\section{Discussion}\label{sec:discussion}

\subsection{Two regimes of commutativity}

Our results reveal a fundamental dichotomy in the role of commutativity
for multi-step proof systems.

\begin{enumerate}[leftmargin=*,itemsep=2pt]
\item \textbf{Error accumulation.}
  Commuting (abelian) error operators yield lower Lyapunov exponents
  (Theorem~\ref{thm:lyapunov}).  This means that when each proof step
  introduces a small perturbation, independent substeps whose errors
  commute produce slower error growth.

\item \textbf{Proof space exploration.}
  Non-abelian groups can have larger spectral gaps
  (Theorem~\ref{thm:spectral_gap}), corresponding to faster mixing of
  the associated random walk.  For proof search, this means that
  non-commuting operations can explore the space of candidate proofs
  more efficiently.
\end{enumerate}

These two objectives---minimizing error accumulation and maximizing
exploration---are fundamentally in tension.  A proof strategy must
balance them: use commuting substeps where errors compound (e.g., in
algebraic simplification chains) and non-commuting operations where
diversity of exploration is needed (e.g., case analysis or induction
strategy selection).

\subsection{Connections to prior work}

\paragraph{Lie group error propagation.}
Theorem~\ref{thm:lyapunov} extends the log-linear error framework
of Barrau and Bonnabel \cite{barrau2014} from continuous-time Lie
group systems to discrete stochastic matrix products.  While
\cite{barrau2014} showed that group-affine systems enjoy linear error
dynamics on the Lie algebra, our result quantifies the Lyapunov
exponent for the discrete setting and shows that the abelian/non-abelian
distinction has concrete consequences for convergence rates.

\paragraph{Quantum error correction.}
The requirement in Gottesman's stabilizer formalism \cite{gottesman1996}
that the stabilizer group be abelian parallels our finding that abelian
error structure yields tighter convergence bounds
(Theorem~\ref{thm:diaconis}).  In both settings, commutativity provides
the algebraic structure needed for systematic error control.
Arvind et al.\ \cite{arvind2002} further showed that abelian subgroups
of the non-abelian error group suffice for constructing error-correcting
codes, analogous to our observation that it is the commutativity of the
error directions (not the full proof transformations) that matters.

\paragraph{Random walks on groups.}
Our Theorem~\ref{thm:spectral_gap} applies the spectral gap framework
of Diaconis \cite{diaconis1988} to the proof search setting.  The
$4{:}1$ asymptotic ratio between the spectral gaps of $D_n$ and
$\Z_{2n}$ illustrates how non-abelian Cayley graphs serve as better
expanders, a phenomenon well-known in the theory of expander graphs
but not previously connected to proof-theoretic convergence.

\paragraph{Error propagation on motion groups.}
Wang and Chirikjian \cite{wang2008} showed that error distributions
on non-abelian Lie groups exhibit coupling between mean and covariance
evolution.  Our commutator-based gap estimate
(Theorem~\ref{thm:lyapunov}\ref{item:gap}) quantifies this coupling
for discrete products and identifies the expected commutator Frobenius
norm as the precise invariant controlling the penalty.

\subsection{Implications for proof systems}

Our results suggest three design principles for iterative proof
refinement systems.

\begin{enumerate}[leftmargin=*,itemsep=2pt]
\item \textbf{Decompose into commuting substeps.}
  When a proof can be decomposed into independent subgoals whose error
  operators commute, error accumulation is minimized
  (Theorem~\ref{thm:lyapunov}\ref{item:diag}).  This favors proof
  strategies based on independent lemmas over tightly coupled reasoning
  chains.

\item \textbf{Use non-commuting operations for search.}
  When the proof strategy involves exploring a space of candidate proofs,
  non-commuting operations (case splits, induction choices) create
  better-connected search graphs (Theorem~\ref{thm:spectral_gap}),
  enabling faster convergence to a correct proof.

\item \textbf{Monitor the commutator norm.}
  The quantity $\E[\frobnorm{\comm{A_1}{A_2}}^2]$ is a computable
  invariant that predicts the convergence penalty from proof step
  interactions (Corollary~\ref{cor:commutator}).  Proof systems could
  estimate this quantity online and adjust their strategy accordingly.
\end{enumerate}

\subsection{Open questions}

We conclude with several directions for future work.

\begin{conjecture}\label{conj:constant}
The dimension-dependent constant $C_d$ in
Theorem~\ref{thm:lyapunov}\ref{item:gap} satisfies
$C_d = \Theta(1)$ as $d \to \infty$, i.e., the Lyapunov gap is
dimension-independent at leading order.  Computational evidence
(Table~\ref{tab:dimension}) supports this.
\end{conjecture}

\begin{table}[ht]
\centering
\caption{Lyapunov exponent gap $\lambda_{\mathrm{gen}} - \lambda_{\mathrm{abel}}$
  as a function of dimension $d$ at $\eps = 0.05$.}\label{tab:dimension}
\begin{tabular}{@{}rc@{}}
\toprule
$d$ & $\lambda_{\mathrm{gen}} - \lambda_{\mathrm{abel}}$ \\
\midrule
2  & 0.000864 \\
4  & 0.001230 \\
6  & 0.001238 \\
8  & 0.001192 \\
10 & 0.001124 \\
\bottomrule
\end{tabular}
\end{table}

\begin{enumerate}[leftmargin=*,itemsep=2pt]
\item \textbf{Discrete proof spaces.}
  Our results use continuous matrix groups.  Extending the framework
  to finite proof state spaces (e.g., propositional proof systems)
  with a purely discrete group structure would make it directly
  applicable to formal verification systems such as Lean or Coq.

\item \textbf{Optimal proof decomposition.}
  Given a proof goal, what is the algorithmically optimal decomposition
  into substeps that minimizes the commutator norm of the resulting
  error operators?  This optimization problem connects our framework
  to proof complexity theory.

\item \textbf{Connection to proof complexity.}
  Can the Lyapunov exponent or spectral gap of the error propagation
  operator be formally related to proof length or proof width in the
  sense of propositional proof complexity?

\item \textbf{Empirical validation.}
  Testing the predictions of Theorems~\ref{thm:lyapunov}
  and~\ref{thm:spectral_gap} on actual proof search traces from
  interactive theorem provers would bridge the gap between our
  theoretical framework and practical proof engineering.
\end{enumerate}
